\documentclass[12pt]{article}
\usepackage{amsmath}
\usepackage{amsfonts}
\usepackage{amssymb}
\usepackage{geometry}
\usepackage{tikz}
\geometry{a4paper, margin=1in}
\pagenumbering{gobble}

\title{02YMECH - Homework 8}
\date{Assigned for the week of Nov 10, 2025}
\author{}
\begin{document}
\maketitle

\section*{Question}
\begin{enumerate}
    \item (Hard) Perform an explicit calculation of time average of potential energy over a one complete period for a particle moving in an elliptical orbit under a central inverse square law force field. Express the result in terms of the force constant of the field and the semi-major axis of the ellipse. Perform a similar calculation for the time average of kinetic energy. Compare your results and verify that they are consistent with the virial theorem.
    \item Two particles are moving under the influence of their mutual gravitational force describe circular orbits about one another with a period $\tau$. If they suddenly stopped in their orbits and allowed to gravitate toward each other, show that they would collide after a time interval $\tau/4\sqrt{2}$.
    \item A thin rod of length $L$ lies along the $x$–axis from $x = 0$ to $x = L$.  
Its linear mass density is not uniform; instead, it varies as a power of the distance from the left end:
\[
\lambda(x) = k\,x^{n}, \qquad 0 \le x \le L,
\]
where $k>0$ is a constant and $n>-1$ is a real parameter.

\begin{enumerate}
    \item Compute the total mass $M$ of the rod in terms of $k$, $n$, and $L$.
    \item Find the position of the center of mass $x_{\mathrm{CM}}$ as a function of $n$ and $L$.
    \item Suppose that the center of mass is required to be located at $x_{\mathrm{CM}} = \frac{3L}{4}$. Determine the value of $n$ for which this holds. Describe qualitatively how the mass is distributed along the rod for this value of $n$.
\end{enumerate}

\end{enumerate}
\end{document}
